% Work/MAT -- 1. rocnik -- hodina 39 -- finalni revize 2026-10-30.
% Priklady z uzivatelem dodanych obrazku upraveny a doplneny vlastnimi ulohami.
% Vysledky a podminky kontrolovany nezavisle; zapis vysledku nepredpoklada
% kladnost promennych tam, kde to neni stanoveno.
\documentclass[11pt,a4paper]{article}
% Spolecny styl pro ucitelske pripravy MAT (Work/MAT)
% Verze 1.1, 2026-10-23
\NeedsTeXFormat{LaTeX2e}
\ProvidesPackage{priprava}[2026/10/23 v1.1 Sablona ucitelskych priprav MAT]

% Zakladni typografie a cestina
\RequirePackage[T1]{fontenc}
\RequirePackage[utf8]{inputenc}
\RequirePackage[czech]{babel}
\RequirePackage{lmodern}
\RequirePackage{microtype}
\RequirePackage{amsmath,amssymb}
\RequirePackage{xcolor}
\RequirePackage{enumitem}
\RequirePackage[most]{tcolorbox}
\RequirePackage{tabularx,booktabs,array}
\RequirePackage{fancyhdr}
\RequirePackage{needspace}
\RequirePackage{geometry}
\RequirePackage[hidelinks]{hyperref}

\geometry{a4paper,top=17mm,bottom=18mm,left=23mm,right=23mm}
\definecolor{notegray}{gray}{0.39}
\setlength{\parindent}{0pt}
\setlength{\parskip}{2.5pt plus 1pt minus 0.5pt}
\setlist[itemize]{topsep=2pt,itemsep=1pt,parsep=0pt,leftmargin=1.45em}
\setlist[enumerate]{topsep=2pt,itemsep=2pt,parsep=0pt,leftmargin=1.65em}
\setcounter{secnumdepth}{0}
\renewcommand{\arraystretch}{1.12}

% Zahlavi: vlevo rocnik | globalni tema, vpravo cislo: nazev hodiny
% Zapati zustava zcela prazdne.
\pagestyle{fancy}
\fancyhf{}
\renewcommand{\headrulewidth}{0.35pt}
\renewcommand{\footrulewidth}{0pt}
\setlength{\headheight}{14pt}
\newcommand{\nastavPripravu}[4]{%
  \fancyhead[L]{\small #1\enspace\textbar\enspace #2}%
  \fancyhead[R]{\small #3: #4}%
}

% Minimalisticke cernobile boxy. Kazda definice zvlast.
\tcbset{commonlesson/.style={enhanced,breakable=false,colback=white,colframe=black,
  colbacktitle=white,coltitle=black,boxrule=0.45pt,
  arc=0.8mm,outer arc=0.8mm,left=3mm,right=3mm,top=1.6mm,bottom=1.6mm,
  toptitle=1mm,bottomtitle=0.7mm,fonttitle=\bfseries,
  attach boxed title to top left={xshift=2.5mm,yshift=-1.15mm},
  boxed title style={colback=white,colframe=white,boxrule=0pt,left=1mm,right=1mm}}}
\newtcolorbox{zadaniBox}[1][]{commonlesson,title={Zadání},#1}
\newtcolorbox{prikladBox}[1][]{commonlesson,title={Příklad},#1}
\newtcolorbox{definiceBox}[1][]{commonlesson,title={Definice},#1}

% Spolecny box pro zadani a jeho reseni, bez mezery mezi castmi.
% Pouziti:
%   \begin{prikladReseniBox}
%     Zadani ulohy.
%     \tcblower
%     \begin{align*} ... \end{align*}
%   \end{prikladReseniBox}
% Spodni cast ma jemny sedy podklad a delici linku; lze v ni sazet
% vice radku matematiky. Obe casti maji jen jeden spolecny ramecek.
\tcbset{solutionlesson/.style={
  commonlesson,
  enhanced jigsaw,
  breakable,
  colbacklower=black!3,
  segmentation style={solid,draw=black!25,line width=0.35pt},
  fontlower=\small\color{notegray},
  before lower={\noindent\textbf{Řešení:}\enspace},
  leftlower=3mm,rightlower=3mm
}}
\newtcolorbox{prikladReseniBox}[1][]{solutionlesson,title={Příklad},#1}
\newtcolorbox{zadaniReseniBox}[1][]{solutionlesson,title={Zadání},#1}

% Pomucky pro autora pripravy: sedive poznamky; \reseni pro starsi pripravy
\newcommand{\poznamka}[1]{\par\begingroup\small\color{notegray}\itshape
  \textnormal{\textbf{Poznámka pro učitele.}}\enspace #1\par\endgroup}
\newcommand{\reseni}[1]{\par\begingroup\small\color{notegray}
  \textbf{Očekávané řešení:}\enspace #1\par\endgroup}
\newcommand{\blok}[1]{\needspace{5\baselineskip}\subsection*{#1}}
\newcommand{\vyrok}[1]{\uv{#1}}

\nastavPripravu{1. ročník}{Mocniny a odmocniny}{39}{Pracovní list}
\begin{document}

\textbf{Jméno a příjmení:}\ \rule{7cm}{0.3pt}\hfill
\textbf{Třída:}\ \rule{2.1cm}{0.3pt}

\textbf{Samostatné procvičování před testem.}
Zjednoduš každý výraz, uveď \textbf{podmínky} původního zadání.
Není nutné všechno spočítat během jedné hodiny: použij vybrané
položky ve škole a zbytek doma. K zápisu výpočtů použij vlastní papír.

\blok{Sada A -- algebraické výrazy a záporné exponenty}
\begin{enumerate}[label=\textbf{\arabic*.},leftmargin=10mm]
\item $\displaystyle\frac{a^2b}{c^3}:\frac{ab^2}{c^2}$.
\item $\displaystyle\frac{a^{-3}b^2}{c^{-3}}:\left(\frac{b^3c}{a^2}\right)^2$.
\item $\displaystyle\left(\frac{9a^{-2}c^3}{8b^{-2}}\right)^{-3}$.
\item $\displaystyle\frac{2a^5b^3}{3a^2b}:\frac{ab}3$.
\item $\displaystyle\frac{(x+y)^3(x+y)^4}{(x+y)^5}:(x+y)^2$.
\end{enumerate}

\blok{Sada B -- mocniny prvočísel místo velkých čísel}
\begin{enumerate}[label=\textbf{\arabic*.},leftmargin=10mm,start=6]
\item $\displaystyle\frac{8^2\cdot6^3}{4^3\cdot12}$.
\item $\displaystyle\frac{125^2\cdot2^{10}}{16^2\cdot50}\cdot\frac{25}{5^3}$.
\item $\displaystyle\left(\frac{2^2}{5}\right)^5\left(-\frac{5^2}{2^3}\right)^3$.
\item $\displaystyle\frac{3^7\cdot9^{-2}}{27^{-1}}$.
\item Výsledek vyjádři \emph{jako součin celočíselných mocnin prvočísel}:
\[
\left(\frac{8\cdot10^2\cdot9^{-3}}{15^4\cdot12^{-3}}\right)^{-3}
:\left(\frac{24^{-1}\cdot6^3}{20^2\cdot5^{-4}}\right)^2.
\]
\end{enumerate}

\blok{Sada C -- odmocniny, úpravy a usměrňování}
\begin{enumerate}[label=\textbf{\arabic*.},leftmargin=10mm,start=11]
\item $\sqrt{200}-2\sqrt{18}+\sqrt8$.
\item $\sqrt[3]{-250}+2\sqrt[3]{16}$.
\item $\displaystyle\frac{\sqrt{75}-\sqrt{12}}{\sqrt3}$.
\item $\displaystyle\frac1{2-\sqrt3}-\frac1{2+\sqrt3}$.
\item Pro $x\in\mathbb R$ uprav $\sqrt{(x-4)^2}$.
\end{enumerate}

\blok{Sada D -- kombinace a technický zápis}
\begin{enumerate}[label=\textbf{\arabic*.},leftmargin=10mm,start=16]
\item $(\sqrt{12}+\sqrt{27})(\sqrt3-1)$.
\item $81^{-3/4}\cdot27^{2/3}$.
\item $\displaystyle\frac{0{,}00084}{2{,}1\cdot10^{-6}}$.
\item Šest souborů po $256\,\mathrm{KiB}$: kolik je to $\mathrm{MiB}$?
\item $\displaystyle\frac{3x^{-2}y^4}{6x^3y^{-1}}$; uveď podmínky.
\end{enumerate}

\newpage
\blok{Výsledky k pracovnímu listu -- samostatná část}
\textbf{Klíč pro kontrolu:} u zlomků podmínky patří k \emph{původnímu}
zápisu, i když se po zjednodušení některý jmenovatel vykrátí.

\textbf{Sada A -- proměnné}
\begin{enumerate}[label=\textbf{\arabic*.},leftmargin=10mm]
\item $\displaystyle\boxed{\frac a{bc}}$, $a,b,c\ne0$.
\item $\displaystyle\boxed{\frac{ac}{b^4}}$, $a,b,c\ne0$.
\item $\displaystyle\boxed{\frac{512a^6}{729b^6c^9}}$, $a,b,c\ne0$.
\item $\boxed{2a^2b}$, $a,b\ne0$.
\item $\boxed1$, $x+y\ne0$.
\end{enumerate}

\textbf{Sada B -- mocniny prvočísel}
\begin{enumerate}[label=\textbf{\arabic*.},leftmargin=10mm,start=6]
\item $\boxed{18}$.
\item $\boxed{250}$.
\item $\boxed{-10}$.
\item $\boxed{3^6=729}$.
\item $\boxed{2^{-25}\cdot3^{17}\cdot5^2}$.
\end{enumerate}

\textbf{Sada C -- odmocniny}
\begin{enumerate}[label=\textbf{\arabic*.},leftmargin=10mm,start=11]
\item $\boxed{6\sqrt2}$.
\item $\boxed{-\sqrt[3]2}$.
\item $\boxed3$.
\item $\boxed{2\sqrt3}$.
\item $\boxed{|x-4|}$ pro všechna $x\in\mathbb R$.
\end{enumerate}

\textbf{Sada D -- kombinované úlohy}
\begin{enumerate}[label=\textbf{\arabic*.},leftmargin=10mm,start=16]
\item $\boxed{15-5\sqrt3}$.
\item $\boxed{1/3}$.
\item $\boxed{400}$.
\item $\boxed{1{,}5\,\mathrm{MiB}}$ (při $1\,\mathrm{MiB}=1024\,\mathrm{KiB}$).
\item $\displaystyle\boxed{\frac{y^5}{2x^5}}$, $x,y\ne0$.
\end{enumerate}

\blok{Krátká kontrola postupu u náročnějších úloh}
\begin{itemize}
\item \textbf{3:} záporný vnější exponent převrátí celý zlomek;
$9^{-3}8^3a^6b^{-6}c^{-9}=512a^6/(729b^6c^9)$.
\item \textbf{10:} převodem na prvočísla vyjde
$\bigl(2^{11}3^{-7}5^{-2}\bigr)^{-3}:
\bigl(2^{-4}3^25^2\bigr)^2=2^{-25}3^{17}5^2$.
\item \textbf{14:} oba jmenovatele mají po násobení sdruženým výrazem
součin $4-3=1$.
\item \textbf{20:} exponentský součet v čitateli je třeba snížit
o exponenty ve jmenovateli; $x^{-2-3}y^{4-(-1)}=x^{-5}y^5$.
\end{itemize}
\end{document}
